\clearpage

\section{设计复盘与结论}\label{sec:7}

\subsection{依据轨迹定位问题}\label{sec:7.1}

同样是未到点，问题可能出在定位、全局规划、局部控制、资源准入或交付验收。以下结合轨迹和事件，按五个层次说明如何确定应修改的部分：
\begin{itemize}
\item \textbf{第一层：位姿与坐标关系}\quad AMCL 需要提供可用的 map—odom 修正，车辆才能在同一张地图上理解目标点、代价地图和传感器扫描。如果这一层不稳定，全局路线即使在地图上正确，也可能在车体坐标下落到错误方向，此时调整 DWB 评价项没有意义。定位结果为后续规划提供共同参考，任务完成仍由验收条件判断。
\item \textbf{第二层：全局路径与代价地图}\quad 静态障碍形成的不可进入栅格、障碍层与同伴层给出的实时占据、膨胀层增加的贴障代价要分别看。若路径贴着障碍，应先区分实际间隙不足还是高代价区域促使规划器改道；直接缩小膨胀半径会减少运动间隙，在窄通道里不能代替几何测量。
\item \textbf{第三层：局部控制与速度采样}\quad 检查局部控制能否沿全局路线选出可执行的速度命令。车原地摆动时，应核对评价项是否引起姿态牵引与路径推进冲突；车几乎不动时，则应分别排查障碍评分、加减速约束和局部目标设置。不同原因都表现为控制器未给出期望速度，却需要修改不同环节。
\item \textbf{第四层：资源状态与入口动作}\quad 车已拿到通道授权却长期滞留在入口附近时，延长预约名义时窗并不能解决进场动作本身的问题，需要结合轨迹判断是否调整对准点以及是否与队友路径产生几何冲突（见第~\ref{sec:3.2} 节）。将授权、目标下发与实际跨口分开记录，才能区分问题属于调度许可还是入口运动。
\item \textbf{第五层：任务验收与服务确认}\quad 车辆到达目标邻域后，仍可能因低速保持不足、驻留未完成或服务对象被占用而尚未交付。同一辆车的通道进入、通道清空和最终交付分别对应资源占用与任务服务两类条件，不应合并为单一“到达时间”，否则会漏掉跨口之后的服务过程。
\end{itemize}

\subsection{从诊断到规则改进}\label{sec:7.2}

\subsubsection*{基线整场规则改进}

最终规则来自三类问题：差速车航向调整异常、通道口反复报占用、观察过期后车辆持续等待。对应的定位依据、修改和复验列在附录~\ref{sec:C}。

早期两次八实体诊断分别完成 11 项和 13 项任务，后一次记录五次通道占用重叠。轨迹显示，车身在参考航点 W、E 附近驻留时可能跨过原有占用窗口边缘。原窗口为 $2.05\le x\le4.95\,\mathrm m$，走廊墙范围为 $1.8\le x\le5.4\,\mathrm m$。正式试验改按车体包络判空，在通道外对准和等待，并采用驻留迟滞与短时脱困。正常回合最终得到 14/14 按期、5 次进入、5 次清空及 0 次占用重叠。

首份观察延迟同为 7.3~s 时，只拒收过期帧会让 B2 一直等到 400~s 窗口结束，按期数为 13/14。开启补采后，B2 在 314.240~s 交付，按期数恢复到 14/14。另一个两车对照中，前车清空即放行的做法比等到固定 32~s 少等 24.112~s。

这几处调整说明，单车动作、共享资源和任务时间必须放在一起检查：单车控制决定车辆如何行驶，通道互斥决定车辆如何让行，任务窗口决定哪些等待可以接受。判据还应按电机实际能达到的精度制定，而不能按理想点位设定。早期回合的偏差帮助定位了问题，第~\ref{sec:4} 章的正式结果再由同条件复跑确认。

\subsubsection*{对抗专项分层排障}

对抗专项的调试遵循同样的分层思路，正式八组回合前共排除三类问题（汇总见附录~\ref{sec:C}）：
\begin{itemize}
\item \textbf{通信读写竞争}\quad 第一批回合中红方无人机在约 40~s 停止响应，原因为以 16~ms 周期轮询观察文件与写入方冲突；给总线读取加共享锁并将轮询放慢至 250~ms 后恢复正常。
\item \textbf{负对照窗口设计}\quad 虚占至 120~s 释放时，无防卫组仍为 14/14，未显露损失；将滞留推迟至 145~s 后无防卫组降为 12/14，防卫效果有了可比较的基准。
\item \textbf{并行仿真环境}\quad 四路并行时处理器争用使基线回归墙钟时间翻倍并引发 R1 超期，改为每批最多两路后时序稳定。
\end{itemize}

\subsection{设计结论与任务效能}\label{sec:7.3}

参照课程中提出的“技术性能、战术（任务）效能、操作（使用）效用”三层次评估度量体系 \cite{ref2}，表~\ref{tab:38} 汇总本次设计在三个层次上的检验结果。

\begingroup\small
\setlength{\tableavailable}{\dimexpr\linewidth-4\tabcolsep\relax}
\begin{longtable}{@{}>{\raggedright\arraybackslash}p{0.20000\tableavailable}>{\raggedright\arraybackslash}p{0.40000\tableavailable}>{\raggedright\arraybackslash}p{0.40000\tableavailable}@{}}
\caption{本次设计检验到的三个层次}\label{tab:38}\\
\toprule
\textbf{评价层次} & \textbf{此次设计中的问题} & \textbf{本次依据} \\
\midrule
\endfirsthead
\multicolumn{3}{@{}l}{\small\color{muted}表~\thetable\quad （续）}\\
\toprule
\textbf{评价层次} & \textbf{此次设计中的问题} & \textbf{本次依据} \\
\midrule
\endhead
\midrule
\multicolumn{3}{r@{}}{\footnotesize\color{muted}续下页}\\
\endfoot
\bottomrule
\endlastfoot
技术性能 & 地面车能否沿路线行驶、停稳、把通道完整让出？ & 传感反馈、轨迹、低箱绕行、运输驻留和各回合清空记录 \\
任务效能 & 单车行动能否合成按时完成的队伍任务？ & 基础回合、低箱与旋翼接管均 14/14 按期；响应阈值防卫在虚占下三轮恢复 14/14 \\
使用效用 & 在任务更多、条件更杂或出现对抗干扰时，流程是否仍有效？ & 起点变化与断链恢复均 14/14，持续断链其余 13 项按期；帧认证拦截伪造帧并维持 14/14 \\
\end{longtable}
\endgroup

\subsubsection*{主要结论}

\noindent\textbf{单体执行与共享通道}\quad 差速控制与双高度雷达配合局部重规划，使地面车能够修正行驶、绕开地图外低箱并完成驻留交付，多旋翼在传感反馈下的悬停误差约 0.053~m；按车体包络判空并依次放行后，三组基础回合及全部补充、对抗回合均为零通道占用重叠，按实际清空放行比固定时段等待少等 24.112~s。

\noindent\textbf{空地观察与对抗防卫}\quad 空中观察采样时刻直接决定蓝方第二批运输何时发车，过期帧拒收后的主动补采使按期任务数由 13/14 恢复到 14/14，持续断链时则将安全等待的代价留在任务分母中。对抗回合表明，通道虚占不引发碰撞却使按期数降至 12/14，采用社会性昆虫响应阈值模型配合裁判闲置核验可在 2—5~s 内收回名额并使三轮随机种子均恢复 14/14；重放欺骗在按期数上不留痕迹，需借助独立审计记录的放行观察真实年龄识别，纳入采样时刻的 HMAC 认证拦截伪造帧且未增加正常延迟。离散算例另表明，岗位支援是否值得实施取决于任务期限与新增车时之间的权衡。

\noindent\textbf{事件审计与综合评价}\quad 在给定场地里，通道占用区间、按期任务数、时间余量与放行观察的真实年龄，分别反映通行效率、队伍任务效能、时序余量与信息可信度。将准入、清空、采信与撤销都落实为可核查事件后，可在同一任务卡上通过调节首帧时延、侦察车起点、入口间隔、资源争用或对抗开关检验不同约束下的协同行为。

\subsubsection*{适用范围与后续方向}

实验结论的适用边界与后续扩展方向如下：
\begin{itemize}
\item \textbf{调度简化}\quad 固定中转点、预设接手车和固定放行顺序简化了调度，动态选车与多路径调序有待在整场中进一步验证；
\item \textbf{软件故障与认证假设}\quad 观察链路故障与重放均为软件注入，未涉及射频传播、定位信号欺骗与密钥管理更新；
\item \textbf{真值核验迁移}\quad 闲置核验依赖仿真真值，迁移到实车需替换为车载定位与里程计联合判断；
\item \textbf{统计与博弈扩展}\quad 除阈值防卫三轮种子与正常条件三轮复跑外，多数对照采用单次回合；后续可在更多随机种子与场地布局下开展多次统计试验，并引入红蓝双方策略自适应的多轮博弈。
\end{itemize}
